\documentclass{article}

\usepackage[margin=1in,top=0.9in,bottom=0.9in]{geometry}
\usepackage{amsmath}
\usepackage{microtype}
\usepackage{booktabs}
\usepackage{enumitem}
\usepackage{titlesec}
\usepackage{parskip}
\usepackage{hyperref}

\titlespacing{\section}{0pt}{8pt}{4pt}
\titlespacing{\subsection}{0pt}{6pt}{2pt}

\setlength{\parskip}{4pt}
\setlength{\parindent}{0pt}

\title{\textbf{Workshop Proposal: Efficient Memory for Long-Context LLM Inference}\\
\large NeurIPS 2026 Workshop}

\author{
  David Aror \quad \texttt{tejirid.aror@gmail.com}\\
  Independent Research
  \and
  Xavistin C. Arul Arasu \quad \texttt{js01.christopher@gmail.com}\\
  Independent Research
  % TODO: add further co-organizers before submission
}

\date{}

\begin{document}
\maketitle

%% ============================================================
%% MAIN PROPOSAL (target: <=3 pages)
%% ============================================================

\section*{1\quad Title and Workshop Description}

\textbf{Title:} \emph{Efficient Memory for Long-Context LLM Inference: Compression, Streaming, and Theoretically-Grounded Methods}

\textbf{Tagline (140 chars):}
\textit{Where exact guarantees meet practical streaming: online KV-cache compression, low-rank methods, and the next 128K-token frontier.}

\subsection*{Topic and Content}

Autoregressive transformer inference at 32K--128K+ token contexts is bottlenecked
by the KV cache, which grows linearly in context length and dominates memory for
deployed models.  A 7B-parameter model requires $\approx$4.7\,GB for the KV cache
alone at $T{=}8{,}192$ tokens in float32---an amount that scales linearly with
context.  As production serving shifts toward much longer contexts, KV cache
management becomes the binding constraint on batch size, GPU utilization, and
serving cost.

The space of solutions has fragmented rapidly.  \emph{Offline/batch} methods
(post-sequence SVD, cross-layer factorizations, quantization-aware low-rank
decomposition) are theoretically well-understood but incompatible with streaming
generation.  \emph{Online/streaming} methods (Oja's rule, eviction policies,
sliding windows, attention sinks) are deployment-compatible but often lack
formal guarantees or achieve compression only at significant quality cost.
\emph{Hardware-aware} approaches (FlashAttention variants, paged attention,
speculative decoding with cache reuse) interact with compression in
underexplored ways.

A critical gap is the absence of a venue where theoreticians who study
low-rank matrix approximation, practitioners deploying at scale, and
systems researchers building inference kernels can exchange ideas in real time.
Recent arXiv activity (OjaKV, SVDq, xKV, KVTC, KQ-SVD, StreamingLLM, SnapKV,
PyramidKV, MagicPIG) reflects a field that is generating results faster
than the community can synthesize them.

\textbf{Why now?} Three developments make 2026 the right moment:
(i) context windows of 128K--1M tokens are now standard in frontier models,
pushing KV memory from a convenience issue to a hard engineering constraint;
(ii) a first wave of principled streaming methods (combining randomized linear
algebra with online updates) has just appeared, and early cross-paper comparison
shows significant quality variation with no agreed-upon benchmark protocol;
(iii) hardware vendors are shipping memory-bandwidth-optimized silicon whose
impact on compression trade-offs is not yet understood.
Without a coordinating venue, the field risks parallel rediscovery, incompatible
benchmarking conventions, and missed connections between theory and practice.

\subsection*{Scope and Audiences}

The workshop targets four overlapping communities:
\begin{itemize}[nosep,leftmargin=1.5em]
  \item \textbf{Theory/algorithms:} randomized linear algebra, streaming algorithms, online learning, matrix sketching.
  \item \textbf{ML systems:} inference engines (vLLM, TensorRT-LLM, TGI), XLA/CUDA kernel authors.
  \item \textbf{LLM practitioners:} fine-tuning, RLHF, long-document QA, multi-turn dialogue, RAG pipelines.
  \item \textbf{Hardware:} accelerator architects interested in memory-hierarchy trade-offs.
\end{itemize}

\textbf{Problems we aim to advance:}
\begin{enumerate}[nosep,leftmargin=1.5em]
  \item Establish shared benchmarking protocols for online KV compression (reconstruction error, end-task perplexity, latency, memory footprint).
  \item Identify which theoretical guarantees (EYM optimality, Halko error bounds, Oja convergence) translate meaningfully to downstream task quality.
  \item Surface open questions at the theory-systems interface: XLA-compilable streaming updates, paged-attention integration, rank adaptation across layers.
\end{enumerate}

\subsection*{Concurrent Submissions}

The lead organizer has not submitted concurrently related proposals to other venues.
% TODO: each organizer must verify and complete this statement.

\subsection*{Location Preference}

\begin{enumerate}[nosep,leftmargin=1.5em]
  \item Sydney
  \item Atlanta
  \item Paris
\end{enumerate}

\subsection*{Invited Speakers}

Organizers will confirm definitive speaker commitments before the proposal
deadline.  The list below reflects initial outreach; confirmed speakers are
marked \textbf{(C)}, tentative \textbf{(T)}.

% TODO: fill in confirmed speakers once commitments are received
\begin{tabular}{lll}
\toprule
Speaker & Affiliation & Status \\
\midrule
Speaker A: online KV / streaming algorithms & TBD & (T) \\
Speaker B: randomized linear algebra        & TBD & (T) \\
Speaker C: ML systems / inference serving   & TBD & (T) \\
Speaker D: long-context LLM benchmarking    & TBD & (T) \\
Speaker E: industry practitioner            & TBD & (T) \\
\bottomrule
\end{tabular}

We target 4--5 invited speakers with at most 25-minute slots to leave ample time
for contributed talks, a poster session, and open discussion.

\subsection*{Diversity and Inclusion}

% TODO: complete once co-organizers are confirmed
The organizing committee will include researchers spanning career stage
(junior PhD students through senior faculty/industry researchers),
gender, geographical region (North America, Europe, Asia-Pacific), and
sector (academia and industry).  We will:
\begin{itemize}[nosep,leftmargin=1.5em]
  \item Issue an explicit diversity statement in the call for papers.
  \item Actively recruit invited speakers from underrepresented groups and regions.
  \item Apply for NeurIPS social-good travel grants to subsidize attendance for researchers from under-resourced institutions.
  \item Use double-blind review for contributed papers.
\end{itemize}

\subsection*{Estimated Attendance}

We estimate 150--300 attendees, based on the intersection of the NeurIPS
efficient-inference, long-context LLM, and randomized algorithms communities.
No previous iteration of this specific workshop has been held.

\subsection*{Special Requirements}

A poster session area for 20--30 posters is requested.
No unusual A/V or compute requirements beyond standard presentation setup.

\subsection*{Differentiation from Related Workshops}

Recent NeurIPS/ICML workshops on efficient transformers (e.g., EfficientML,
sparsity-focused workshops) have emphasized pruning, quantization, and
distillation.  KV cache compression, particularly streaming/online methods with
theoretical guarantees, has not been the primary focus of any recent workshop.
Our workshop is distinguished by: (i) exclusive focus on memory---not speed,
sparsity, or model compression broadly; (ii) explicit emphasis on bridging
theory (randomized NLA, online learning) with practice (inference engines,
benchmarks); (iii) a live benchmarking session where submitted methods are
evaluated on a shared held-out protocol announced at submission time.

\subsection*{Preliminary Benchmark Findings (Pilot for the Live Benchmark Session)}

To validate that the shared protocol in item~(1) above is measurable and
discriminative, the organizers ran a pilot sweep on a single commodity
workstation using GPT-2-small key/value activations extracted from WikiText-2.
The pilot compares two streaming compression methods head-to-head:
\textbf{TKV}, a novel warmup-then-freeze approach that computes an
EYM-optimal SVD basis during prefill and freezes it for the decode phase, and
\textbf{OjaKV}~\cite{ojakov2025}, the current state-of-the-art online baseline
based on Oja's streaming rule.

The key insight behind TKV is that autoregressive decode tokens are
generated \emph{conditioned on} the prefill context and therefore tend to remain
within the subspace already spanned by prefill activations.  Freezing the
SVD basis after warmup exploits this structure: the prefill establishes an
EYM-optimal low-rank subspace; subsequent tokens are projected onto it in
$O(Rd)$ per token rather than re-running a streaming SVD update.  This avoids
OjaKV's stale-basis failure mode, in which the rotating Oja basis invalidates
previously computed projections as new tokens arrive.

The pilot exercises every axis of the proposed protocol---reconstruction error,
memory footprint, end-task perplexity, and per-token latency---and confirms
that these axes separate the two methods cleanly:

\begin{itemize}[nosep,leftmargin=1.5em]
  \item \textbf{Per-token decode latency.}  After a shared warmup of 256 tokens,
    TKV processes each new decode token in $0.17$--$0.18$\,ms (rank $8$--$32$),
    versus $0.45$--$0.56$\,ms for OjaKV---a $2.5$--$3.2\times$ speedup that is
    \emph{flat across ranks} because the frozen projection cost is $O(Rd)$
    independent of the SVD update schedule.
  \item \textbf{Reconstruction error (real activations).}  OjaKV's relative
    Frobenius error exceeds $1.2$ at every rank tested (rank $8$--$32$,
    $T{=}512$\,tokens of decode), meaning its reconstruction is worse than the
    zero matrix---direct evidence of stale-basis collapse.  TKV's error
    is $0.27$--$0.44$ over the same range, $3$--$4.5\times$ lower.  At long
    context ($T{=}32$K, rank $16$) TKV achieves $0.55$ versus OjaKV's $0.63$.
  \item \textbf{Memory at scale.}  At $T{=}128$K tokens the full float32 KV cache
    requires $9.44$\,GB; a TKV pyramid basis compresses this to $2.21$\,GB
    ($4.3\times$) while beating OjaKV on reconstruction error at every long-context
    length tested ($1$K--$32$K).
  \item \textbf{End-task perplexity (GPT-2).}  TKV (average rank $15$
    pyramid) achieves perplexity $89.6$ versus $129.6$ for the exact cache (ratio
    $0.69$)---the frozen basis acts as a mild regularizer and \emph{improves}
    held-out perplexity.  OjaKV under the same pyramid schedule achieves $756.2$
    (ratio $5.83\times$ worse than exact), confirming that stale-basis error
    compounds into catastrophic language-model degradation.
  \item \textbf{Long-context perplexity (TinyLlama-1.1B, WikiText-2).}
    In a sliding-window evaluation at context lengths $512$--$2048$ tokens,
    TKV perplexity ranges from $36$ to $658$ while OjaKV degrades to
    $932$--$1795$, consistent with the stale-basis collapse observed above.
\end{itemize}

These numbers are a single-workstation pilot, not a leaderboard; their purpose is
to demonstrate that the live-benchmark session can rank submitted methods on a
shared, reproducible protocol.  Full scripts and figures accompany this proposal.

\subsection*{Workshop Schedule (Proposed)}

\begin{tabular}{ll}
\toprule
Time & Session \\
\midrule
08:30--09:00 & Opening remarks and problem framing \\
09:00--11:00 & Invited talks (2 $\times$ 25 min + discussion) \\
11:00--12:00 & Contributed talks (3 $\times$ 15 min) \\
12:00--13:30 & Lunch + Poster session \\
13:30--15:00 & Invited talks (2 $\times$ 25 min + discussion) \\
15:00--16:00 & Live benchmark reveal and community discussion \\
16:00--17:00 & Panel: ``What does theory owe practice, and vice versa?'' \\
17:00--17:30 & Contributed talks (2 $\times$ 12 min) \\
17:30--18:00 & Open discussion and closing \\
\bottomrule
\end{tabular}

Contributed work will be solicited via OpenReview (required per NeurIPS guidelines).
All accepted submissions will be non-archival; authors may optionally post on arXiv.
Accept/reject notification will be issued by \textbf{September 29, 2026} at the latest.

%% ============================================================
%% ORGANIZER INFORMATION (target: <=2 pages)
%% ============================================================

\newpage
\section*{Organizer Information}

\subsection*{Organizing Committee}

\textbf{David Aror} \quad \texttt{tejirid.aror@gmail.com} \quad Independent Research

David Aror is an independent researcher working on efficient inference for large
language models, with a current focus on online KV cache compression.  His recent
work introduces \textbf{TKV}, a streaming KV compression method that computes
an EYM-optimal low-rank basis during prefill (via Brand's rank-1 SVD update and
Halko randomized range finding) and freezes the basis for the autoregressive decode
phase---exploiting the structure of conditional generation to achieve $2.5$--$3.2\times$
lower per-token latency than OjaKV with $3$--$4.5\times$ lower reconstruction error.
He has experience running small research communities and organizing reading groups
on efficient ML.  He is not submitting any other workshop proposals to NeurIPS 2026.

\textbf{Xavistin C. Arul Arasu} \quad \texttt{js01.christopher@gmail.com} \quad Independent Research

Xavistin C. Arul Arasu is an independent researcher working on efficient inference
systems for large language models, with a focus on streaming KV-cache compression
and the empirical evaluation of low-rank and sink/window hybrid methods.  He led the
benchmarking effort behind this proposal---reconstruction-error, memory-ratio,
perplexity, and streaming-error sweeps comparing exact streaming SVD (Brand update)
against Oja's-rule baselines on GPT-2---which forms the basis for the shared
live-benchmark protocol described in Section~1.  He is not submitting any other
workshop proposals to NeurIPS 2026.

% TODO: add further co-organizers with complementary expertise (systems, theory, industry)
% Format: Name, affiliation, email, paragraph on expertise and organizing experience.
% Ensure team covers: gender diversity, geographic diversity, career stage diversity.
% Confirm none are on more than 2 total workshop proposals.

\subsection*{Program Committee}

The PC will be recruited to ensure each submission receives at least 3 reviews and
no reviewer is assigned more than 3 papers.  We estimate 60--100 submissions
(based on submission rates at adjacent workshops) and therefore need 60--100
reviewers.

% TODO: list confirmed PC members with affiliations
Confirmed PC members (to be updated before submission):
\begin{itemize}[nosep,leftmargin=1.5em]
  \item PC Member 1, Affiliation \textbf{(C)}
  \item PC Member 2, Affiliation \textbf{(C)}
  \item \ldots\ (target: 60+ members to cover expected submission volume)
\end{itemize}

Organizers will recuse themselves from reviewing any submission with which they
have a conflict of interest as defined by the NeurIPS CoI policy.
Specifically, no organizer will be involved in assessing submissions from their
own organization.  Conflict detection will use OpenReview's built-in CoI tooling.

Workshop organizers will not submit papers to this workshop (per NeurIPS policy).

%% ============================================================
\bibliographystyle{plain}
\begin{thebibliography}{9}

\bibitem{ojakov2025}
Anonymous.
OjaKV: Online KV cache compression via Oja's rule.
arXiv:2509.21623, 2025.

\bibitem{halko2011}
N.~Halko, P.-G.~Martinsson, J.~A.~Tropp.
Finding structure with randomness.
\emph{SIAM Review}, 53(2):217--288, 2011.

\bibitem{brand2006}
M.~Brand.
Fast low-rank modifications of the thin SVD.
\emph{Linear Algebra and its Applications}, 415(1):20--30, 2006.

\bibitem{streamingllm2023}
G.~Xiao et al.
Efficient streaming language models with attention sinks.
arXiv:2309.17453, 2023.

\bibitem{svdq2025}
Anonymous. SVDq. arXiv:2502.15304, 2025.

\bibitem{xkv2025}
Anonymous. xKV. arXiv:2503.18893, 2025.

\end{thebibliography}

\end{document}

